\documentclass[a4paper,unpublished,onecolumn,11pt]{quantumarticle}
\pdfoutput=1

\usepackage{amsmath}
\usepackage{amssymb}
\usepackage{amsthm}
\usepackage[normalem]{ulem}
\usepackage{comment}

\usepackage{env}

\newcommand{\Pauli}{\mathcal P}

\title{Verification of universal quantum computations through single-qubit tomography}
\author{Cool people}
%\author{Ulysse Chabaud}
%\affiliation{}
%\author{Dominik Leichtle}
%\affiliation{}

\begin{document}


\begin{abstract}
    The abstract.
\end{abstract}


\maketitle

\begin{comment}

\section{Summary of the project}

\begin{itemize}
    \item A verification of quantum circuits using Pauli measurements.
    \item The circuits we can verify are composed of any tensor product input (and magic states in particular), followed by any combination of $CZ$ gates.
    \item The measurements needed for the verification are randomly chosen $X_i\otimes Z^{\otimes\backslash i}$ (X measurement of the $i^{th}$ qubit and $Z$ measurements of all the other qubits).
    \item Alternatively, the number of measurement settings can be reduced up to two, depending on the chromatic number of the graph corresponding to the CZ gates.
    \item The idea of the verification is as follows: we combine straightforward single-mode fidelity estimation using $X$ and $Z$ measurements with a lemma yielding a fidelity witness for tensor product target states. Then, we reduce the general verification to that of the input tensor product state by inverting the entangling operation at the level of classical post-processing (since the eigenspaces of the operators of the form are conserved by products of $CZ$ gates, in a way which is efficiently computable).
    \item We get polynomial security for the individual state preparation, but we can boost the security to exponential if these states are used either for \textsf{BQP} computations or specific sampling tasks.
    \item We should obtain a verification of quantum speedup, and, adding adaptive measurements, a blind (since the resource state should be universal) verification of universal quantum computations with magic states.
    \item Strikingly, our verification protocol is not very different from the original receive-and-measure protocol, but works when the input $\ket+$ states are replaced by arbitrary single-qubit states.
    \item The verification method not only gives an accept/reject answer but is actually robust (it gives a tight and robust lower bound on the fidelity with the target state).
\end{itemize}

\begin{figure}[h!]
    \centering
    \includegraphics[width=0.7\columnwidth]{MS verification.pdf}
    %\caption{Caption}
    %\label{fig:MSverif}
\end{figure}

\section{What we can look at}

\begin{itemize}
    \item Check the existing literature.
    \item Check the power of these circuits, including the case of adaptive measurements.
    \item Write the security proof, including the many-copy case using De Finetti-type theorems. What is the best fidelity estimation technique with $X$ and $Z$ measurements out there?
    \item How important is the structure of the graph? It seems that a bipartite graph would allow a much more efficient verification since the number of $X$ and $Z$ measurements for each qubit would be balanced (which is probably better for tomography).
    \item Which tests (i.e. measurement settings) and distribution of tests gives the best detection probability? Do we achieve the inverse of the fractional chromatic number of the graph, as for prepare-and-send verification with standard traps? Can we even go beyond that (compare to dummyless testing strategies)?
    \item Can we extend these results to more general circuits? Like generic Clifford circuits instead of only $CZ$ gates? Also, what happens for $CCZ$ gates? Should work for circuits that are $(Cliff)^{\otimes n}-CZ-(Cliff')^{\otimes n}$.
        \item Include $Y$ measurements as well? For that, check that the corresponding eigenstates are also mappedt to Pauli eigenstates in a balanced way.
    \item Can we relate these circuits to physical Hamiltonians/platforms? 
    \item Similarly, could there be a nice experimental demonstration of these techniques? This protocol should be particularly useful since it is the same setup used for tomography.
    \item Can we think of other applications than speedup demonstration and universal computing? Possibly using different input states we could get states that are useful for something else, such as distributed quantum computing, quantum communication, sensing, or secret sharing?
    \item More applications: can we use fidelity with such target states as witnesses for specific quantum properties, like entanglement, $T$ count, or negativity of the discrete Wigner function?
    \item The works that seem to achieve a similar result all deal with hypergraph states. How do we compare to these and what are our advantages and disadvantages compared to them?
    \item Can the protocol be upgraded to fault-tolerance with an equivalent version in logical space?
\end{itemize}

\end{comment}
\section{Alternative notes}
\subsection{What do we do}
\begin{itemize}
    \item We certify almost any quantum state $\ket\Psi$, specifically states of the form $\ket\Psi = C(\bigotimes_{i\leq n} \ket{\psi_i})$ where $C$ is a Clifford unitary on $n$ qubits using $O(n)$ single-qubit Pauli measurements.
    This means any $n-$qubit state of that form can be \textit{efficiently} certified using tools of \textit{single-qubit tomography}.
    \item The same procedure allows to verify universal quantum computations. Indeed, the model "product state + Clifford gate + single-qubit measurements with classical post-processing of the outcomes" is universal. So the same task we ask from the server (preparing $C(\bigotimes_{i\leq}\ket{\psi_i})$ -- untrusted), complemented by the task we can do at home (measuring in a Pauli basis of our choice, therefore possibly the result of a classical-post processing of previous measurements outcomes -- trusted), is enough to perform any quantum computation.
    
    \item Therefore, because we are able to certify the state prepared by the server (which is the only untrusted entity), we can \textit{efficiently verify any quantum computation through single-qubit tomography}.

    \item Is it good ? 
    \begin{itemize}
        \item The task of certifying a quantum state is important. We'd like to do it through single-qubit measurements, \textit{i.e} the simplest quantum setting to certify a quantum state.
        \item Direct Fidelity Estimation often gives a sample complexity that is exponential in $n$ in the worst case \cite{FL11:direct, SLP11:practical, FOS24:optimal}, with \cite{FOS24:optimal} stating an exponential lower bound for Haar-random states.
        \item Other methods exist, and are based on, instead of giving an estimate of the fidelity, they give a property that relates to the fidelity. In \cite{HPS24:certifying}, they call it a \textit{shadow overlap}, and are used to certify \textit{almost} any $n-$qubit quantum state with a sample complexity $O(n^2)$ (states that relate to a quantum Markov chain with relatively slow mixing time). In \cite{CGKM21:efficient, UC24:efficient}, they call it a \textit{robust fidelity witness} and they use it to efficiently certify bosonic systems. In this paper, we use the same robust fidelity witness technique generalized to the $n-$qubit case and show that it allows to efficiently certify \textit{almost} all $n-$qubit quantum states, namely any product state that has gone through a $n-$qubit Clifford unitary, using $O(n)$ samples.
    \end{itemize}
    \item How do we do it concretely ?
    \begin{itemize}
        \item We know that Direct Fidelity Estimate (DFE) are powerful tools but they are useful in small dimensions only. Here, we use them separately to get $n$ fidelity estimates, that we use to compute a single (robust) fidelity witness for the target state $C(\bigotimes_{i\leq n} \ket{\psi_i})$
        \item Indeed, using DFE on each qubit of $\bigotimes_{i\leq n} \ket{\psi_i}$ separately, to later combine the single-qubit fidelities to build a witness for the target state $\bigotimes_{i\leq n} \ket{\psi_i}$ is straightforward and requires only few single-qubit Pauli measurements.
        \item We show that we can use the same setup and same number of samples to certify $C(\bigotimes_{i\leq n} \ket{\psi_i})$, because each single-qubit Pauli measurement in the previous setting can be mapped in a tensor product of Pauli measurements in the current setting, using the conjugation of the Pauli observable by the Clifford circuit $C$.
    \end{itemize}
    \item Optimization:
    \begin{itemize}
        \item Some measurements can be performed simultaneously, so we also provide a procedure that aims to minimize the total number of required samples.
        \item In this optimization procedure, we show how we very simply obtain a very powerful result, namely that the number of samples depends only on the required precision of the fidelity (the typical Hoeffding-style result), but the number of sample does not depend on the size of the computation!!! this is a huge result (apparently).
    \end{itemize}
\end{itemize}
\subsection{Notations and preliminaries}
\begin{itemize}
    \item Single-qubit Pauli $\rightarrow$ Pauli string of weight $1$: For $P\in \Pauli_1$
    \begin{equation}
        P_i = 
        \left({\bigotimes_{j<i} \mathds{1}}\right)
        \otimes P \otimes
        \left({\bigotimes_{j>i} \mathds{1}}\right)
    \end{equation}
    \item The outcome $x$ of the measurement described by the Pauli string $\prod_i P_i$ is the product of the different outcomes $\{x_i\}_i$ obtained when measuring $P_i$ separately:
    \begin{equation}
        x = \prod_i x_i
    \end{equation}
\item Characteristic function associated to a $d-$dimensional state $\rho$ : a function of Pauli operators
\begin{equation}
   \begin{array}{rccl}
     \chi_\rho  :&\Pauli_n   &\rightarrow& \mathbb Z \\
                &P          &\mapsto    &  \Tr(\rho P)/\sqrt{d} \\
    \end{array} 
\end{equation}
\item Single-qubit fidelity estimator: let $N_i$ be the number of samples required to estimate the fidelity of qubit $i\leq n$. Then, let $j$ denote one of the sample's number. Draw the single-qubit Pauli $P^{(j)}$ from $\mathcal{D}_i$ (defined below), and let $P_i^{(j)}$ denote its Pauli-string version. The random variable $y_{ij}$ corresponds to the outcome of measuring $P_i^{(j)}$ on the $j-$th sample of the state. Here, $\chi_{\rho_i}(\cdot)$ is the characteristic function associated to the single-qubit state $\rho_i = \Tr_{\backslash i} [\rho]$.
   \begin{equation}
   \label{eq: estimator}
        \bar X_i = \dfrac{1}{N_i}\sum_{j=1}^N y_{ij}\; \mathrm{sgn}(\chi_{\rho_i}(P ^{(j)})) \; ||\chi_{\rho_i}||_1 
    \end{equation}
\item Normalization factor
    \begin{equation}    
    \label{eq: normalization}
    ||\chi_{\rho_i}||_1 = \sum_{P\in \Pauli_1} |\Tr[\rho_i P]|
    \end{equation}
\item Probability distribution associated to the single-qubit state $\rho_i$
\begin{equation}
    \label{eq: state prob. distribution}
            \mathcal D_i(P) = \dfrac{
        |\chi_{\rho_i}(P)| 
    }{
        ||\chi_{\rho_i}||_1
    }
        \end{equation}
 \item Fidelity witness (definition):
 \begin{equation}
 \label{eq: fidelity witness}
     W(\rho, \sigma) = 1- \sum_{i=1}^n (1 - F(\rho_i, \sigma_i))
 \end{equation}
 \item Fidelity witness (empirical):
 \begin{equation}
 \label{eq: fidelity witness empirical}
     \bar W(\rho, \sigma) = 1- \sum_{i=1}^n (1 - \bar X_i)
 \end{equation}
 \item Support of a Pauli operator $P\in \Pauli_n$
 \begin{equation}
 \label{eq: support of Pauli}
     i\in S(P) \iff [P]_i\neq \mathds{1}
 \end{equation}
\end{itemize}
\subsection{(Naive) Protocol for product states}
Here, $\ket\psi = \bigotimes_{i=1}^n \ket{\psi_i}$ is the target state, and $\rho$ is its density matrix. $\sigma$ is a state sent by the Server to the Client. The protocol goes as follows, with fixed Client inputs being the target state, and the precision $\varepsilon, \delta$.\\

\begin{itemize}
    \item For $i\leq n$:
        \begin{itemize}
            \item $\ket{\psi_i}$ is the target state. The Client sets $N_i = ||\chi_{\rho_i}||_1^2 /\varepsilon^2 \; \log(1/\delta)$, where $\rho_i=\ket{\psi_i}\bra{\psi_i}$, and computes the normalization factor $||\chi_{\rho_i}||_1$ and the probability distribution $\mathcal D_i$ associated to the state, according to \ref{eq: normalization} and \ref{eq: state prob. distribution}. Then, for $j\leq N_i$:
            \begin{itemize}
                \item The Client samples $P_{i}^{(j)}$ according to $\mathcal D_i$, and asks for a copy $\sigma$.
                \item The Client measures the received state $\sigma$ with the operator $P_i^{(j)}$ and gets outcome $y_{ij}$.
            \end{itemize}
            \item The Client computes the estimator $\bar X_i$ following \ref{eq: estimator}. This serves as an estimate of $F(\rho_i, \sigma_i)$.
        \end{itemize}
    \item Then, the Client computes $W(\rho, \sigma)$ following \ref{eq: fidelity witness}.
\end{itemize}

Analysis of the protocol:
\begin{itemize}
    \item Cost of the protocol: number of copies of $\sigma$ asked
    \begin{equation}
        N = \sum_{i=1}^n N_i
    \end{equation}
    Note: this is a very naïve version of the protocol because using one copy of $\sigma$, we measure only one qubit, and throw away the rest. If, instead of that, using the fact that we are dealing with a product state, we notice that one can do the sampling beforehand, and measure qubits simultaneously, this can reduce the overall cost. Concretely, for $i=1, ..., n$, for the $j-$th copy of $\sigma$ requested, one can sample $\{P_{i}^{(j)}\}_i$ in advance, compute the string $P^{(j)} = \prod_i P_i^{(j)}$ and perform the $n$ measurements on the same copy. Hence, the number of copies required will be given by $\max_i(N_i)$.
    
    % For a product state and using one copy to perform $n$ measurements instead of a single one, we can get this to $N = \max_i(N_i)$ pretty easily.
    \item Assumptions of the protocol: i.i.d assumption, to get rid of using De Finetti reduction
    \item Security of the protocol: $\bar X_i$ is $\varepsilon-$close to $F(\rho_i, \sigma_i)$ with probability $1-\delta$, proved using Hoeffding's inequality. Then, [to complete]
\end{itemize}

\subsection{Protocol for any Clifford circuit}
Let $C$ be a $n-$qubit unitary from the Clifford group. Now, we assume the target state is $C(\ket \psi)$: it is not necessarily a product state anymore. We could uncompute $C$ on $\sigma$ and evaluate its fidelity to $\rho$, the target product state, but the Client doesn't have the capability of performing such an operation. Instead, the Client will apply this Clifford transformation on the Pauli measurements, that remain Pauli measurements by definition of the Clifford group. Is this called the Heisenberg picture ?
% In the following protocol, let 
\begin{enumerate}
    \item For $i\leq n$:
        \begin{enumerate}
            \item $\ket{\psi_i}$ is the target state. The Client asks $N_i$ copies of $\rho$,  computes the normalization factor $||\chi_{\rho_i}||_1$ (for $\rho_i = \ket{\psi_i}\bra{\psi_i}$) and the probability distribution $\mathcal D_i$ associated to the state, according to \ref{eq: normalization} and \ref{eq: state prob. distribution}. Then, for $j\leq N_i$:
            \begin{itemize}
                \item The Client samples $P_{i}^{(j)}$ according to $\mathcal D_i$, computes $\tilde P_i^{(j)} = C P_{i}^{(j)} C^\dagger$, and asks for a copy $\sigma$.
                \item The Client measures the received state $\sigma$ with the operator $\tilde P_i^{(j)}$, gets the collection of outcomes $\{y_{ijk}\; : k\in S(\tilde P_i^{(j)})\}$, and sets $y_{ij} = \prod_{k\in S(P_i^{(j)})} y_{ijk}$.
            \end{itemize}
            \item The Client computes the estimator $\bar X_i$ following \ref{eq: estimator}. This serves as an estimate of $F(\rho_i, \sigma_i)$.
        \end{enumerate}
    \item Then, the Client computes $\bar W(\rho, \sigma)$ following \ref{eq: fidelity witness}. This gives an estimate for $W(\rho, \sigma)$.
\end{enumerate}
State a theorem about this protocol, relating the $N_i$ to the precision and overall error probability.

\begin{theorem}[Efficiency of Protocol]
    test
\end{theorem}
\subsection{Merging Pauli strings for fewer measurements}
Preliminaries: 
\begin{itemize}
    \item \textit{Compatibility condition of Pauli operators}
    \item \textit{Compatible Pauli operators describe compatible measurements}, \textit{i.e} several measurements that can be performed simultaneously using single-qubit Pauli measurements.
\end{itemize}

Let $C_G$ now be a Clifford circuit composed of $CZ$ gates only, realizing an entanglement of the input qubits according to a graph $G$ that is $k-$colorable. This has the following consequences:
\begin{itemize}
    \item For sampled Pauli measurements $Z_i$ and $Z_j$, their back-propagated version $\tilde Z_i$ and $\tilde Z_j$ are compatible measurements, because $Z$ operators commute with $CZ$. This can be generalized: any collection of measurements $\{Z_i\}$ can be performed simultaneously if $C=C_G$ for a given $G$. 

    \item For $X_i$ and $X_j$, $\tilde X_i$ and $\tilde X_j$ are compatible if and only if $(i,j)\notin G$. This is because $X$ doesn't exactly commute with $CZ$, its commutation makes a $Z$ appear on the neighboring qubit : $CZ_{ij} X_i = X_i Z_j CZ_{ij}$. Therefore $\tilde X_i = X_i Z_{N(i)}$, where $N(i)$ denotes the set of neighboring qubits of $i$. As $\tilde X_j = X_j Z_{N(j)}$, we conclude that $\tilde X_i$ and $\tilde X_j$ are compatible iff $j\notin N(i)$ and vice-versa.
    \item A direct consequence of the previous observation is that any given collection of $X-$measurements $\{X_i\}_{i\in c}$ can be performed simultaneously if $C = C_G$ and $c$ is a color of $G$.
\end{itemize}

We use the previous facts to draw the following protocol.

\subsection{To-do}
\begin{itemize}
    \item De Finetti reduction to remove i.i.d assumption
    \item Fuchs-van de Graaf, and link fidelity estimation to the security of the computation
    \item Nuance between overlap and fidelity estimation
    \item What does "back-propagating observable" really mean ? Is it really a $\oplus$ of the outcomes?
    \item Some figures!
\end{itemize}

\subsection{Cryptogrpahic scenario -- removing the i.i.d assumption}
\begin{itemize}
    \item De Finetti: will not work because scales poorly with the local dimension that is here $d=2^n$. The old beginning of the notes:
\begin{itemize}
    \item Generalization of de Finetti's representation theorem for discrete variables quantum systems \cite{Renner_2007}
    \item The theorem works on permutation-invariant systems. Problem: the state given by the Server is not necessarily permutation-invariant! so to make it fit in de Finetti's assumptions, we do the following trick: independently of the Server's sent states, we choose among the $N$ rounds which ones will be used for test. This \textbf{enforces} the permutation-invariance. We randomly shuffle the $N$ rounds, and we choose a random subset of $k$ rounds: the reduced state $\rho_k = \Tr_{N-k}[\tilde \rho_N]$ is permutation-invariant; 
    where $\tilde \rho_N$ is a mixture among all the possible permutations:
    \begin{equation}
        \tilde \rho_N = \dfrac{1}{N!}\sum_{\pi\in S_N} 
        \Pi_\pi
        \rho_N
        \Pi_\pi^\dagger
    \end{equation}
    \item Indeed, state $\tilde \rho_N$ is permutation-invariant, so any subsystem also inherits from this property. Namely: Partial traces of symmetric states are symmetric.
\end{itemize}
\item Alternative solution: use Serfling bound on classical random variables. It says the following, for random variables (such as measurement outcomes, or a bounded function of the measurement outcomes) in the range $[a, b]$. If you have a population of size $N$, that you randomly divide in subsets of size $n$ (the tested subset) and $m$ (untested) with $N=n+m$, then with probability $1-\delta$ you get
\begin{equation}
    |
    \bar X_{rem}
    -
    \bar X_{n}
    |
    \leq
    \dfrac{N}{m}(b-a)
    \sqrt{
    \dfrac{N-n+1}{2nN}\log{\dfrac{2}{\delta}}
    }
\end{equation}
where $\bar X_n$ is the average obtained on the empirical data, and $\bar X_{rem}$ is the average on the untested samples.
\item It requires to know the range of the random variable, requires it to be bounded
\item Main question: what does it tell us on a single element of the untestet set, if it was sampled uniformly at random ? (the rest being discarded)
\end{itemize}

\newpage
\section{Certification of product states through single-qubit tomography}
Essentially using single-qubit fidelity witnesses in DV. Also can be seen as the particular case $C=\mathbb I$ of the general protocol presented later.
\begin{itemize}
    \item \textbf{Problem}: we desire to prepare a target state $\ket\Psi = \bigotimes_{i<n}\ket{\psi_i}$. With a source giving an unknown $n-$qubit quantum state $\rho$, what procedure can we do to estimate $F(\rho, \ket \Psi)$, using only Pauli measurements ? Typically this task requires several copies of the state $\rho$ to estimate $F(\rho, \ket \Psi)$ up to precision $\varepsilon$ with probability $\delta$. A relevant metric is thus the number of copies of $\rho$ required to achieve such a task. 
    % We want a procedure to build an estimator $\bar X$ of $F(\rho, \ket \Psi)$ using $N(\varepsilon, \delta)$ copies of $\rho$ and allowed to perform Pauli measurements on them, and such that
    % \begin{equation}
    % \label{eq: Estimator precision}
    %     \mathrm{Pr}\left[
    %         |\bar X - F(\rho, \ket \Psi)| \geq \varepsilon
    %     \right] \leq \delta
    % \end{equation}
    \item \textbf{Computing a fidelity witness through single-qubit fidelities}. The task of estimating $F(\rho, \ket\Psi)$ can be \textbf{reduced} to the estimation of $F(\rho_i, \ket{\psi_i})$ for $i<n$ then recombining the results in a specific procedure presented in \cite{CGKM21:efficient} to build a witness of $F(\rho, \ket\Psi)$, meaning a number $W(\rho, \ket\Psi)$ that lower bounds the fidelity, and that is lower-bounded by a linear function of the fidelity:
    \begin{equation}
        f(F) \leq W(\rho, \ket \Psi) \leq F
    \end{equation}
    with
    \begin{equation}
        W(\rho, \ket\Psi) = 1- \sum_{i=1}^n (1-F(\rho_i, \ket{\psi_i}))
    \end{equation}
    and $f(F) = 1-n(1-F(\rho, \ket\Psi))$. Therefore, the task of estimating the fidelity of a multi-qubit product state can be reduced to that of estimating the single-qubit fidelities separately.
    \item \textbf{Benefits from the reduction.}   
    Procedures to estimate the fidelity of any $n-$qubit state to any $n-$qubit target state exist in the literature like in \cite{FL11:direct, FOS24:optimal}, with the most notable result being that the fidelity could be estimated in a number of copies of $\rho$ that scales exponentially in $n$. Therefore, using these methods to estimate the fidelity of large systems might be impractical, while they would be very efficient for small systems, like to estimate single-qubit fidelities.
    
\end{itemize}

\subsection{Estimate single-qubit fidelities}
Protocol to estimate $F(\rho_i, \ket{\psi_i})$. Let us first observe that, with $\sigma_i=\ket{\psi_i}\bra{\psi_i}$
\begin{equation}
    F(\rho_i, \ket{\psi_i}) = \sum_{P\in\Pauli_1} \Tr[\rho_iP_i]\Tr[\sigma_iP_i]
\end{equation}
where $P_i$ is a Pauli string of size $n$ where the only non-identity Pauli is $P$ at index $i$. 
We have full knowledge of the state $\sigma$ but $\rho$ is sent by the server so we can compute $\Tr[\rho_iP_i]$ empirically only. By slightly modifying this expression, this task can be reduced to computing an empirical average of a random variable. Indeed, let $||\chi_{\sigma_i}||_1 = \sum_{P\in \Pauli_1} |\Tr[\sigma_i P_i]|$. Then, the fidelity can be re-written, following the approach of \cite{FOS24:optimal}, as 
\begin{equation}
    F(\rho_i, \ket{\psi_i}) = \underset{P\sim \mathcal H}{\mathbb E}[g(P)]
    % \sum_{P\in\Pauli_1} 
    % p(P)\;
    % g(P)
    %
\end{equation}
where the probability distribution $\mathcal H$ is defined as
\begin{equation}
    \mathcal H(P)=
    \dfrac{
        |\chi_{\sigma_i}(P_i)| 
    }{
        ||\chi_{\sigma_i}||_1
    }
    \leq 1
\end{equation}
and the function $g$ is defined as
\begin{equation}
    \label{eq: black box fct}
    g(P) =
        \mathrm{sgn}(\chi_{\sigma_i}(P_i)) \; ||\chi_{\sigma_i}||_1 
    \Tr[\rho_iP_i]\; .
\end{equation}
One can see that introducing $||\chi_{\sigma_i}||_1$ is to yield a valid probability distribution. As $\rho$ is available only experimentally, one can only perform the measurement $P_i$ a total number of $N$ times to approximate $g(P)$ for a given $P$.
Now let there be $N\in \mathbb N$. For $j\leq N$, do the following :
\begin{itemize}
    \item Sample $P^{(j)}$ according to $\mathcal H$, denote $P_i^{(j)}$ the corresponding Pauli string.
    \item Let $Y_{P_j^{(i)}}$ be the binary random variable with possible outcomes $\{-1, +1\}$ corresponding to performing measurement $P_i^{(j)}$ on state $\rho$. We use the notation $Y_{ij}$ to make it shorter. Finally, we write $y_{ij}$ be a sample of $Y_{ij}$.
    \item Then, build the estimator
    \begin{equation}
        \bar X = \dfrac{1}{N}\sum_{j=1}^N y_{ij}\; \mathrm{sgn}(\chi_{\sigma_i}(P_i ^{(j)})) \; ||\chi_{\sigma_i}||_1 
    \end{equation}
    \item By Hoeffding's theorem, for $\varepsilon>0$ and $\delta >0$, if
    \begin{equation}
        N= \mathcal O\left(
        \dfrac{||\chi_{\sigma_i}||_1^2 }{\varepsilon^2} \log\left(\dfrac{1}{\delta}\right)
        \right)
    \end{equation}
    then
    \begin{equation}
        \mathrm{Pr}\left[
        \left|{\bar X - \underset{P\sim \mathcal H}{\mathbb E}[g(P)]}\right| \geq \varepsilon
        \right] \leq \delta
    \end{equation}
    \item In particular, as $||\chi_{\sigma_i}||_1 \leq \sqrt{d}=\sqrt 2$, then $N = \mathcal O(1/\varepsilon^2 \; \log(1/\delta))$ samples are needed to estimate $F(\rho_i, \ket{\psi_i})$. But this is the worst case scenario. In general, we keep $N$ depending on $\sigma_i$, meaning that each target state needs a different number of copies to be certified.
\end{itemize}

\subsection{Building the witness}
% \section{Certification of any multi-qubit state}
\section{Verification of universal quantum computations}
\begin{itemize}
    \item \textbf{Universality} is reached with Clifford unitaries and adaptive Pauli measurements, acting on a product state. Let us focus on the first part and certify a $C(\ket\Psi)$ state, as the Client is assumed to be honest and can therefore perform the Pauli measurements as well as adapt them to perform the computation. Universal quantum computations can thus be efficiently verified if $F(\rho, \ket\Psi)$ can be estimated precisely enough with overwhelming probability with a sample complexity that scales polynomially with $n$.
\end{itemize}
% Now, $C\neq \mathbb I$ but is still restricted to the Clifford group. The target state is thus not a product state anymore, 

\newpage
Old version
\section{Verification of multi-qubit states through fidelity witnesses using Pauli measurements}
\begin{itemize}
    \item (Problem)
    We want to verify the preparation of a state of which we have a classical description $\ket{C}$. This method tackles the class of states that can be written as $C\left(\bigotimes_{i=1}^n\ket{\psi_i}\right)$, where $C$ is a $n-$qubit unitary operator from the Clifford group, and $\ket{\psi_i}$ are single-qubit pure states. We note $\ket\psi = \bigotimes_{i=1}^n\ket{\psi_i}$
    \item (Universality)
    This class of states is not restrictive at all: in fact, as the pure states can correspond to any single-qubit state, in particular the $T\ket +$ magic state,
    any $m-$qubit state can be prepared from $m+t=n$ pure states where $t$ of them are magic, using $t$ Pauli measurements and feed-forward.
    \item (Verification)
    In this work, we tackle the task of estimating the fidelity of a given state $\rho$ to the target state :
    \begin{equation}
        F(\rho, C(\ket\psi))
    \end{equation}
    \item (Fidelity estimation using single-qubit fidelity witnesses)
    This task is not trivial at first glance, so rather we show how to achieve it from estimating fidelities of the system qubit by qubit and recombining the results. In short, estimating $F(\rho_i,\ket{\psi'_i})$ using single-qubit measurements, for $\ket{\psi'_i}$ that are efficiently relatable to the initial $\ket{\psi_i}$. For instance, if $C=\mathbb I$, it is known that an estimation of $F(\rho_i, \ket{\psi_i})$ for all $i$ leads to a robust witness of $F(\rho, \ket\psi)$, by the work of [Chabaud 2020]. Indeed, if the collection of single-qubit fidelities are obtained, the following metric can be computed
    \begin{equation}
        \label{eq:fidelity witness}
        W(\rho, \ket\psi) = 1- \sum_{i=1}^n [1-F(\rho_i, \ket{\psi_i}]
    \end{equation}
    and it satisfies the following properties:
    \begin{enumerate}
        \item Robustness property $1-n(1-F(\rho, \ket\psi)) \leq W(\rho, \ket\psi)$: if the fidelity is $1$, the witness is $1$.
        \item Witness property: $W(\rho, \ket\psi) \leq F(\rho, \ket\psi)$
    \end{enumerate}
    They are combined in short in
    \begin{equation}
        \label{eq: robust fidelity witness}
        f(F) \leq W(\rho, \ket\psi) \leq F
    \end{equation}
    where $f$ is a continuous and linear function of $F$ such that $f(1)=1$.

    
    \item (Measurement back-propagation) When $C\neq \mathbb I$, the work of [Upreti, 2024] shows a turnaround: 
    computing the single-qubit fidelities $F(\rho_i, \ket{\psi'_i})$ can be done in the same way as computing $F(\rho_i, \ket{\psi_i})$, simply by classically post-processing the measurement outcomes. This is efficient for instance if the measurements are projective measurements in a Pauli basis and if $C$ is Clifford.


    \item Conditions: the measurements must be tomographically complete. Pauli measurements are such measurements. Indeed, any quantum state can be expended in the Pauli basis:
    \begin{equation}
        \rho = \sum_{\vec{P}\in \mathcal{P}_n} \alpha_{\vec P} \underbrace{P_1\otimes P_2\otimes ...\otimes P_n}_{\vec P}
    \end{equation}
    where the coefficients are $\alpha_{\vec P} = \mathrm{Tr}[\rho \vec P]$, which can be experimentally obtained by measuring ${\vec P}$ on $\rho$.
\end{itemize}

\subsection{Sketch of the protocol}
%(Estimating single-qubit fidelities through single-qubit measurements and efficient numerical computations)

We first assume $C=\mathbb I$.
How to verify a multi-qubit state $\rho$ using single-qubit measurements ? This problem is reduced in estimating $F(\rho_i, \ket{\psi_i})$ for all $i$, using that in \eqref{eq:fidelity witness} to obtain a witness for $F(\rho, \ket \psi)$ that satisfies \eqref{eq: robust fidelity witness}.

\begin{itemize}
    \item Let us evaluate analytically the expression of $F(\rho_i, \ket{\psi_i})$ to expose how it can later be computed using single-qubit measurements. The first step is to express $\rho_i = \Tr_{\backslash i}[\rho]$.
    \item Using the fact that any operator can be decomposed in the Pauli basis with appropriate coefficients, we can write 
    \begin{equation}
        \rho = \sum_{\vec{P}\in \mathcal{P}_n} \alpha_{\vec P} \underbrace{P_1\otimes P_2\otimes ...\otimes P_n}_{\vec P}
    \end{equation}
    where $\alpha_{\vec P} = \mathrm{Tr}[\rho \vec P]$.

    \item Then, to compute $\rho_i$, we need to take the partial trace on indices $j\neq i$.


    When tracing $\rho$ over all indices except $i$, the trace of Pauli operators in $\vec P$ for $j\neq i$ is taken in product, as $\Tr[P_1\otimes P_2] = \Tr[P_1]\Tr[P_2]$. On the other hand, the trace of Pauli operators is zero except for the identity.
    Therefore the only terms in the sum that remain are the Pauli operators that have the identity on indices $j\neq i$.
    \begin{eqnarray}
        \rho_i &=& \Tr_{\backslash i}\left[ 
        \sum_{\vec{P}\in \mathcal{P}_n} \alpha_{\vec P} {P_1\otimes P_2\otimes ...\otimes P_n}
        \right]
        \\
        &=& 
        \sum_{\vec{P}\in \mathcal{P}_n}
        \alpha_{\vec P}
        \Tr_{\backslash i}\left[ 
        {P_1\otimes P_2\otimes ...\otimes P_n}
        \right]
        \\
        &=& 
        \sum_{\vec{P}\in \mathcal{P}_n}
        \alpha_{\vec P} P_i
        \quad 
        \left(\prod_{j\neq i}\Tr[P_j]\right)
        \\
        &=&
        \sum_{P\in \mathcal P_1} \alpha_{P_i}\quad  \underbrace{\mathbb I\otimes ... \otimes P \otimes ...\otimes \mathbb I}_{P_i}
        \\
    \end{eqnarray}
    We use this to re-define $P_i$ as the Pauli operator on $n-$qubits such that the $i-$th component is $P$ and all the others are the identity. In this, the coefficients are $\alpha_{P_i}=\Tr[\rho P_i]$: obtainable by single-qubit measurements on the given state $\rho$.

    \item The single-qubit fidelities $F(\rho_i,\ket{\psi_i})$ can be computed using single-qubit Pauli measurements. Indeed:

    \begin{eqnarray}
        F(\rho_i,\ket{\psi_i}) 
        &=& \Tr[\rho_i \ket{\psi_i}\bra{\psi_i}] \\
        &=& \Tr\left[\sum_{P\in \mathcal P_1} \alpha_{P_i}  P_i \ket{\psi_i}\bra{\psi_i}\right] \\
        &=& 
        \sum_{P\in \mathcal P_1}
        \alpha_{P_i}
        \Tr\left[
        P_i \ket{\psi_i}\bra{\psi_i}\right] \\
        &=& 
        \sum_{P\in \mathcal P_1}
        \alpha_{P_i}
       \bra{\psi_i}P_i\ket{\psi_i}\\
        &=& 
        \sum_{P\in \mathcal P_1}
        \Tr[\rho P_i]
       \bra{\psi_i}P_i\ket{\psi_i}
    \end{eqnarray}

    In each term of the sum, the first factor is obtained experimentally by performing single-qubit Pauli measurements, while the other one is computable numerically on public data, such as the single-qubit Pauli and the classical description of the $i-$th component of the target pure state.
\end{itemize}



In summary, the protocol can be described as:
\begin{enumerate}
    \item Receive $\rho$
    \item Perform single-qubit Pauli measurements $P_i$, compute averages $\Tr[\rho P_i]$.
    \item Numerically, compute $\bra{\psi_i}P_i\ket{\psi_i}$
    \item Compute single-qubit fidelity estimate 
    \begin{equation}
    F(\rho_i,\ket{\psi_i}) 
    =
        \sum_{P\in \mathcal P_1}
    \Tr[\rho P_i]
   \bra{\psi_i}P_i\ket{\psi_i}
    \end{equation}
    \item Compute $W(\rho, \ket \psi)$ using \eqref{eq:fidelity witness}
\end{enumerate}

(Integrating measurement back-propagation)
Now we handle the case where $C\neq \mathbb I$, using measurement back-propagation: when $C$ is from the Clifford group, a similar approach can be conducted, simply by adapting the single-qubit Pauli measurements accordingly. Indeed, the task becomes to estimating $F(\rho, C(\ket\psi)$ which, if we want to relate the given state to a target product state, could re-write $F(C\rho C^\dagger, \ket \psi)$. Therefore, the same approach can be conducted mathematically with $\rho'=C \rho C^\dagger$:
\begin{equation}
    F(\rho'_i,\ket{\psi_i}) 
    =
        \sum_{P\in \mathcal P_1}
    \Tr[\rho' P_i]
   \bra{\psi_i}P_i\ket{\psi_i}
    \end{equation}
The difference is in how the $
    \Tr[\rho' P_i]$ are obtained experimentally, as the client receives $\rho$, not $\rho'=C\rho C^\dagger$, and is not assumed to have the facilities to perform this transformation coherently on $\rho$. Here [to be continued]
\subsection{Next steps}
\begin{itemize}
    \item What about multi-qubit fidelity estimates ? 
    \item What about their overlap, and the link with Pauli noise learning? 
\end{itemize}
\subsection{Questions}
\begin{itemize}
    \item What is concretely to prove in the sketch of the first approach ? Can the validity of \eqref{eq: robust fidelity witness} be taken from [Chabaud 2020] ?
\end{itemize}

\newpage
\section{Notes}
\begin{itemize}
    \item Fidelity estimation through Pauli measurements, already in [FL11], for a target pure state $\rho=\ket{\psi} \bra\psi$ and a received state $\sigma$: average over all Pauli operators of the products of characteristic functions
    \begin{equation}
        F(\rho, \sigma) = \frac{1}{2^n} \sum_{P\in \mathcal P_n} \chi_\rho(P)\chi_\sigma(P)
    \end{equation}
    \item Characteristic function of a $n-$qubit state $\rho$ : function of Pauli operator
    \begin{equation}
        \chi_\rho(P) = \Tr[P\rho]
    \end{equation}
    \item Flammia's approach does exactly that without the measurement back-propagation.
    \item FL11: For stabilizer states, fidelity can be efficiently estimated by performing Pauli measurements on a few copies of $\sigma$, "independently of the system size"
\end{itemize}

\bibliography{bibliography.bib}

\bibliographystyle{abbrv}
\end{document}
